% ============================================================
%  ch1_introduction.tex
%  Chapter 1: Introduction
%  Master's Thesis -- SQL MATCH_RECOGNIZE on Pandas DataFrames
% ============================================================
% !TEX root = Thesis.tex
\chapter{Introduction}
\label{chap:intro}
Row-pattern recognition searches an ordered series of table rows for
relationships that span several events.  It can detect shapes such as V, W,
and U, as well as trends and repeated behavior.  These patterns support
analysis in many domains.  Examples include money-laundering
transaction chains~\cite{petkovic2022fraud}, failures in IoT sensor data, and
user journeys through e-commerce clickstreams.

This chapter starts from the idea itself.  Section~\ref{sec:intro-rowpatterns}
explains what a row pattern is, what the V, W, and U shapes mean, and how the
existing families of languages express such patterns.
Section~\ref{sec:intro-problem} then states the problem this thesis
addresses.


% ----------------------------------------------------------------
\section{Row Patterns in Ordered Data}
\label{sec:intro-rowpatterns}
% ----------------------------------------------------------------
A row pattern describes a shape that appears across several consecutive
rows.  Building one takes three steps.

First, the rows are divided into independent groups --- one customer, one
sensor, one bank card --- and each group is put in order, usually by time.
Second, each kind of row that matters is given a name, and that name is
given a meaning by a condition.  A row whose price is lower than the row
before it can be called \textsc{down}; a row whose price is higher can be
called \textsc{up}.  Third, a pattern says in which order those names may
appear.

The third step is the same idea used to search text.  A search for a word
followed by any digits does not look for one fixed string; it describes a
shape that many different strings can take.  A row pattern does this for
rows instead of characters~\cite{winand_match_recognize,zemke2007pattern}.
Table~\ref{tab:intro-rowlabels} shows the idea on six days of one
customer's prices.

\begin{table}[H]
\centering
\caption{Six ordered rows, each labelled by a condition.}
\label{tab:intro-rowlabels}
{\small
\begin{tabular}{clll}
\toprule
\textbf{Day} & \textbf{Price} & \textbf{Label} & \textbf{Why} \\
\midrule
1 & 100 & \textsc{start} & the row the shape starts from \\
2 &  80 & \textsc{down}  & lower than day 1 \\
3 &  60 & \textsc{down}  & lower than day 2 \\
4 &  55 & \textsc{down}  & lower than day 3 \\
5 &  70 & \textsc{up}    & higher than day 4 \\
6 &  90 & \textsc{up}    & higher than day 5 \\
\bottomrule
\end{tabular}
}
\end{table}

\noindent
The price falls for three days and then rises for two.  In words: one
\textsc{start} row, then one or more \textsc{down} rows, then one or more
\textsc{up} rows.  That sentence is written as a pattern like this:

\begin{lstlisting}[language=SQL]
PATTERN (START DOWN+ UP+)
\end{lstlisting}

Each name in the line is a label from Table~\ref{tab:intro-rowlabels}, and
the order of the names is the order the rows must appear in.  The symbol
\texttt{+} after a name means \emph{one or more rows of that kind}.  So
\texttt{START} matches day~1, \texttt{DOWN+} matches days~2 to~4, and
\texttt{UP+} matches days~5 and~6.

Because \texttt{+} does not fix a number, the same line also matches a fall
lasting two days or ten.  It describes the shape rather than one exact
length, which is what makes it a pattern.  This is the V of
Figure~\ref{fig:intro-shapes}.  Writing \texttt{DOWN\{3\}} instead would
demand exactly three falling days, and \texttt{DOWN?} would make the fall
optional.

The difference from an ordinary filter is the unit of interest.  A filter
asks a question about one row: is this price below ten?  A row pattern asks
a question about a run of rows: does the price fall three times and then
recover?  No single row answers the second question.  The answer lives in
the order and in the relationship between neighbors, so the information
disappears the moment the rows are treated as an unordered
set~\cite{melton2017sqlrpr}.

\subsection{The Shape Vocabulary}
\label{subsec:intro-shapes}
Analysts often name these patterns by the outline they trace when the
measured value is plotted against time.  Three names recur, and
Figure~\ref{fig:intro-shapes} shows them side by side.

\begin{figure}[htbp]
  \centering
  \resizebox{\textwidth}{!}{%
  \begin{tikzpicture}[
      x=1cm, y=1cm,
      axisline/.style={draw=black!45, line width=0.6pt,
                       -{Stealth[length=1.7mm]}},
      curve/.style={draw=blue!55!black, line width=1.15pt,
                    line join=round},
      trough/.style={circle, fill=blue!55!black, inner sep=0pt,
                     minimum size=2.6pt},
      guide/.style={draw=black!30, line width=0.5pt, densely dotted},
      shapename/.style={font=\small\bfseries, text=black!85},
      caption/.style={font=\scriptsize, text=black!70, align=center,
                      text width=3.3cm}
    ]

    % ---------- panel 1: V ----------
    \begin{scope}[shift={(0,0)}]
      \draw[axisline] (0,0) -- (3.3,0);
      \draw[axisline] (0,0) -- (0,2.15);
      \node[font=\scriptsize, text=black!60, anchor=north] at (1.65,-0.08) {time};
      \node[font=\scriptsize, text=black!60, rotate=90, anchor=south]
           at (-0.16,1.05) {value};
      \draw[guide] (1.5,0) -- (1.5,0.35);
      \draw[curve] (0.15,1.75) -- (0.75,1.15) -- (1.15,0.62)
                -- (1.5,0.35) -- (1.85,0.62) -- (2.35,1.2) -- (3.05,1.8);
      \node[trough] at (1.5,0.35) {};
      \node[shapename] at (1.65,-0.78) {V shape};
      \node[caption]  at (1.65,-1.42)
           {one low point:\\ a dip and a recovery};
    \end{scope}

    % ---------- panel 2: W ----------
    \begin{scope}[shift={(4.35,0)}]
      \draw[axisline] (0,0) -- (3.3,0);
      \draw[axisline] (0,0) -- (0,2.15);
      \node[font=\scriptsize, text=black!60, anchor=north] at (1.65,-0.08) {time};
      \draw[guide] (0.85,0) -- (0.85,0.42);
      \draw[guide] (2.35,0) -- (2.35,0.42);
      \draw[curve] (0.15,1.8) -- (0.55,0.95) -- (0.85,0.42)
                -- (1.2,0.95) -- (1.6,1.3) -- (2.0,0.95)
                -- (2.35,0.42) -- (2.7,1.0) -- (3.05,1.8);
      \node[trough] at (0.85,0.42) {};
      \node[trough] at (2.35,0.42) {};
      \node[shapename] at (1.65,-0.78) {W shape};
      \node[caption]  at (1.65,-1.42)
           {two low points:\\ the level is tested twice};
    \end{scope}

    % ---------- panel 3: U ----------
    \begin{scope}[shift={(8.7,0)}]
      \draw[axisline] (0,0) -- (3.3,0);
      \draw[axisline] (0,0) -- (0,2.15);
      \node[font=\scriptsize, text=black!60, anchor=north] at (1.65,-0.08) {time};
      \draw[guide] (0.95,0) -- (0.95,0.42);
      \draw[guide] (2.2,0) -- (2.2,0.42);
      \draw[curve] (0.15,1.8) -- (0.6,0.95) -- (0.95,0.42)
                -- (1.55,0.38) -- (2.2,0.42) -- (2.6,0.98) -- (3.05,1.8);
      \node[trough] at (0.95,0.42) {};
      \node[trough] at (2.2,0.42) {};
      \node[shapename] at (1.65,-0.78) {U shape};
      \node[caption]  at (1.65,-1.42)
           {a flat trough:\\ the low level persists};
    \end{scope}
  \end{tikzpicture}
  }
  \caption{The V, W, and U shapes.}
  \label{fig:intro-shapes}
\end{figure}

\begin{description}
  \item[V shape.] The value falls to a single low point and then rises
    again.  It marks a dip followed by a recovery: a share price that drops
    and rebounds, a temperature that falls and returns, or a customer whose
    spending declines for a few months and then picks up.  The V is the
    smallest interesting shape, because it already needs two opposite
    trends and the turning point between them.

  \item[W shape.] Two V shapes in sequence.  The value falls, recovers only
    partly, falls a second time, and then rises.  In market analysis this
    double bottom is read as a stronger signal than one dip, since the
    level was tested twice.  In monitoring it often means a fault that was
    not really fixed the first time.

  \item[U shape.] The value falls, stays low for a while, and then rises.
    The trough is flat rather than sharp.  The distinction from a V matters
    in practice: a V is a brief disturbance, whereas a U is a sustained
    condition such as an outage, a seasonal lull, or a slow recovery.
\end{description}

The shapes matter because they carry meaning that summary statistics lose.
An average price hides whether the series dipped once, twice, or stayed
low.  Detecting the shape, and reporting where it started and ended, is
what turns an ordered table into an explanation of what
happened~\cite{petkovic2022rpr}.  The same reasoning extends beyond these
three outlines to trends, repeated behavior, and recovery after a failure.

\subsection{Ways to Express a Sequential Pattern}
\label{subsec:intro-approaches}
Sequential patterns are not new, and several families of languages already
describe them.  They differ in how directly the pattern can be stated, and
in what has to be deployed to run it.

\emph{Procedural code} keeps the data in the host language.  An analyst
writes a loop, or shifts columns and combines Boolean masks, and inspects
the rows directly.  Nothing new has to be installed, and any condition can
be expressed.  The cost is that the pattern is never written down as such:
it is spread across index arithmetic, bounds checks, and overlap handling,
where a single off-by-one silently changes the result~\cite{Python_UDF}.

\emph{Window functions and self-joins} stay inside SQL.  A query can reach
the previous or next row with \texttt{LAG} and \texttt{LEAD}, or join a
table to itself once per step.  This is declarative and portable.  But the
pattern must be re-encoded as fixed offsets, so the query grows with the
length of the sequence, a variable-length repetition is awkward to express,
and the plan may scan the table several
times~\cite{cao2012optimizationanalyticwindowfunctions}.

\emph{Complex-event-processing languages} were designed for this problem in
streams.  Systems such as Apache Flink provide pattern operators over
unbounded event sequences~\cite{flink}.  They are expressive, but they
assume a streaming runtime and a separate deployment, which is a large step
for an analyst holding a finished table.

\emph{Row-pattern recognition in SQL} takes the regular-expression idea and
applies it to rows.  The pattern is written once, as a pattern, next to the
data it queries.  SQL:2016 standardized this as the
\texttt{MATCH\_RECOGNIZE} clause~\cite{iso2016,ISO2021}, and it is the
approach this thesis follows.  Consider a card that a thief first tests with
small charges and then empties.  A fraud analyst writes that sequence
directly as the shape it is:

\begin{lstlisting}[language=SQL]
PATTERN (PROBE+ (TRANSFER | WITHDRAW){2,5} DRAIN)
\end{lstlisting}

It is built the same way as \texttt{PATTERN (START DOWN+ UP+)} above, with
two additions.  \texttt{PROBE+} again means one or more test charges.  The
bar in \texttt{(TRANSFER | WITHDRAW)} means \emph{or}, so that row may be
either kind, and \texttt{\{2,5\}} allows between two and five of them.
\texttt{DRAIN} matches the final large movement.  As in Table~\ref{tab:intro-rowlabels}, the names are not
keywords of the language: the analyst invents them and gives each one a
meaning through a condition, so \texttt{PROBE} might mean an amount below a
small threshold.  Chapter~\ref{chap:background} presents the full
clause and its sub-clauses.

The cost of this approach is scope: the clause is optional in the standard,
and the systems that implement it do so to different degrees, which
Chapter~\ref{chap:related} examines in detail.

% ----------------------------------------------------------------
\section{Problem Statement}
\label{sec:intro-problem}
% ----------------------------------------------------------------
SQL:2016 standardized row-pattern recognition through the
\texttt{MATCH\_RECOGNIZE} clause~\cite{iso2016,ISO2021}.  The clause lets a
query define and search patterns within ordered rows.  Oracle had already
introduced it in Oracle Database~12c in 2013~\cite{OracleMR,
oracle2021dwhsg,book}.
Support nevertheless remains uneven.

Trino~\cite{trino2022,trino2021row},
Snowflake~\cite{snowflake2025match_recognize}, and Apache
Flink~\cite{flink} provide their own implementations.  However, support for
the optional SQL:2016 features R010, R020, and R030 varies, as do practical
restrictions.  PostgreSQL~18 lists these features as unsupported, while
DuckDB lists \texttt{MATCH\_RECOGNIZE} as planned work~\cite{
postgresql18unsupported,duckdb2026roadmap}.

This gap matters because many analytical workflows begin in Python.  Users
prepare data, test hypotheses, and build models before moving to a larger
database or cloud platform.  For data that already fits in memory, a separate
SQL engine adds deployment, integration, and data-transfer work.  This thesis
does not aim to replace an enterprise database.  Its main objective is to
bring a practical subset of declarative row-pattern recognition to the Pandas
environment where these workflows often begin.  Section~\ref{sec:intro-objectives}
states that objective precisely and breaks it into three parts.

Without this declarative primitive, SQL users must combine window functions,
self-joins, or loops outside the database.  Beyond two or three steps, the
logic accumulates row-number calculations, timestamp checks, and
\texttt{JOIN} conditions.  It is difficult to maintain and may require
repeated table scans~\cite{cao2012optimizationanalyticwindowfunctions}.

The same limitation appears in Python.  Pandas~\cite{McKinney2010} is a
popular tabular-data library, but it has no built-in SQL row-pattern
primitive.  Analysts instead create shifted columns such as
\texttt{df['value'].shift(-1)} and \texttt{shift(-2)}, combine Boolean masks,
use \texttt{groupby().apply()}, and then pivot or rename intermediate
columns.  Explicit loops are also common.  Listing~\ref{lst:ch1-udf} shows
this approach for a three-row price decrease (\texttt{A~B~C}).

\begin{lstlisting}[language=Python,
    caption={Detecting a three-step price-decrease pattern in
             Pandas \emph{without} \texttt{MATCH\_RECOGNIZE}},
    label={lst:ch1-udf}]
import pandas as pd

df['order_date'] = pd.to_datetime(df['order_date'])
df = df.sort_values(['customer_id', 'order_date']).reset_index(drop=True)

def match_pattern_ABC(group: pd.DataFrame):
    matches = []
    for i in range(len(group) - 2):
        A, B, C = group.iloc[i], group.iloc[i+1], group.iloc[i+2]
        if (B['price'] < A['price']) and (C['price'] < B['price']):
            matches.append({
                'customer_id':  A['customer_id'],
                'start_date':   A['order_date'],
                'end_date':     C['order_date'],
                'bottom_price': C['price']
            })
    return matches

all_matches = []
for _, grp in df.groupby('customer_id'):
    all_matches.extend(match_pattern_ABC(grp))

result_df = pd.DataFrame(all_matches)
\end{lstlisting}

Even this fixed three-row pattern needs several procedural steps.  The
hand-written logic is vulnerable to alignment errors and does not extend
cleanly to longer patterns or patterns with more conditions.  Such cases also
require custom indexing, overlap handling, and output logic.

The same query expressed with \texttt{MATCH\_RECOGNIZE} is shown
in Listing~\ref{lst:ch1-mr}: a single declarative statement that
specifies the partition, ordering, pattern grammar, variable definitions, and
output measures.

\begin{lstlisting}[language=Python,
    caption={The same pattern expressed with \texttt{MATCH\_RECOGNIZE}
             on a Pandas DataFrame},
    label={lst:ch1-mr}]
import pandas as pd
from pandas_match_recognize import match_recognize

df['order_date'] = pd.to_datetime(df['order_date'])

query = """
SELECT customer_id, start_date, end_date, bottom_price
FROM orders
MATCH_RECOGNIZE (
  PARTITION BY customer_id
  ORDER BY     order_date
  MEASURES
    A.order_date AS start_date,
    C.order_date AS end_date,
    C.price      AS bottom_price
  ONE ROW PER MATCH
  AFTER MATCH SKIP TO NEXT ROW
  PATTERN (A B C)
  DEFINE
    B AS B.price < A.price,
    C AS C.price < B.price
)
"""

output_df = match_recognize(query, df)
\end{lstlisting}

The declarative form is not only shorter.  It separates what to find from how
the engine finds it.  The query is easier to audit, compose, and extend.  A
four- or five-step sequence with conditional alternation normally changes the
\texttt{PATTERN} and \texttt{DEFINE} clauses, not the surrounding Python
traversal code~\cite{michels2018new,petkovic2022rpr}.  The implementation can
then compile the pattern into automata, cache repeated pattern structures,
and process ordered partitions without exposing those details to the user.
Chapter~\ref{chap:background} introduces the full structure and
sub-clauses of \texttt{MATCH\_RECOGNIZE}.
% ----------------------------------------------------------------
\section{Research Objectives and Contributions}
\label{sec:intro-objectives}
% ----------------------------------------------------------------
\paragraph{Main objective.}
To bring a practical R010-style subset of SQL:2016 row-pattern recognition to
Pandas, so that a declarative pattern query runs directly on a DataFrame
without a database server.

This connects relational pattern queries with in-memory analytics.  It also
asks a question that the rest of the thesis answers: whether such queries can
run in process, correctly and fast enough to be useful.  Three specific
objectives follow from it.

\paragraph{Specific objectives.}
Three objectives follow.  Each one is stated so that the thesis can be
judged against it.

\begin{enumerate}
  \item \textbf{Determine whether the subset can run in process.}
    The standard defines row-pattern recognition for a database engine, which
    owns its storage, planner, and execution operators.  A DataFrame runtime
    owns none of these.  The first objective is to design and realize an
    engine that parses a \texttt{MATCH\_RECOGNIZE} query, compiles its pattern
    into automata, and executes it over an ordered DataFrame with no database
    server, and so to show whether the standard's execution model survives
    that move.

  \item \textbf{Define the supported subset and verify its correctness.}
    A partial implementation is only useful if its edges are known.  The
    second objective is to state the evaluated R010-style subset precisely
    --- the main clauses, output and skip modes, navigation and aggregate
    functions, and the pattern constructs from quantifiers and alternation to
    exclusions and \texttt{PERMUTE} --- and then to establish that its results
    agree with independent production implementations, rather than with the
    engine's own expectations.

  \item \textbf{Characterize the cost and the limits.}
    Feasibility alone does not justify the approach.  The third objective is
    to measure execution time, throughput, and memory against established
    engines under equal resource limits, and to test how the engine scales as
    the input grows.
\end{enumerate}

Objective~1 is addressed by the design in Chapter~\ref{chap:methodology} and
the implementation in Chapter~\ref{chap:implementation}.  Objectives~2
and~3 are answered by the five evaluation criteria of
Chapter~\ref{chap:evaluation}: correctness and feature coverage for the
second, and usability, performance, and scalability for the third.

\paragraph{Contributions.}
Together, these objectives produce three contributions: an in-process
Pandas engine with tested row-pattern semantics, an optimized and guarded
execution design, and a theoretical and empirical evaluation.
Section~\ref{sec:related-contribution} explains these contributions and
positions them against related systems.  Chapter~\ref{chap:methodology}
describes the engine's five core stages.

\paragraph{Evidence at a glance.}
At the time of final validation, the regression suite contained 709 passing
test cases~\cite{monierashraf2025tests}.  The cases cover pattern families,
constructs, output modes, navigation, and aggregation behavior.  The
controlled cross-system benchmark adds direct comparison with Trino~473 and
Oracle~21c EE.  It uses five pattern families and six nested prefixes of the
Amazon UK Products dataset, from 100\,K rows to the full 2{,}222{,}742 rows.
All 90 system--pattern--size cells completed.  For each of the 30
pattern--size combinations, all three systems returned the same normalized
result set.  Under the shared one-core, 32\,GB memory limit, mean time was
0.17\,s for the engine, 1.04\,s for Oracle, and 1.98\,s for Trino.
Engine time grew almost linearly.  Its measured query-time memory was higher
than either database, because the compiled path builds whole-column masks and
intermediate arrays instead of evaluating row by row.  Its total footprint was
nevertheless lower, because it runs as a single Python process rather than a
resident server that holds the table in a JVM heap or an SGA.

The larger Amazon Reviews test extends the range to 227.9~million rows,
about $100\times$ the main dataset.  Under a shared one-core, 58\,GB
memory limit, the engine and Oracle~21c Enterprise Edition completed all
35 cases.  Trino~473 completed 33.  Its complex nested cases at 160 and
227.9~million rows reached the 600\,s limit.  Engine time remained close to linear.
Four pattern families were completed by all three systems at every size.
Averaging those four at each size, Oracle's mean time was
$5.1$--$10.1\times$ the engine's and Trino's was $11.4$--$13.8\times$ the
engine's.  Oracle's multiple grew with the input, while Trino's stayed
roughly flat.  At the largest size the
four-family means were 28.9\,s for the engine, 291.7\,s for Oracle, and
375.6\,s for Trino.
Chapter~\ref{chap:evaluation} reports the complete method, trade-offs, and
limits.

% ----------------------------------------------------------------
\section{Organization of the Thesis}
\label{sec:thesis-organization}
% ----------------------------------------------------------------

The rest of the thesis is organized as follows.

\textbf{Chapter~\ref{chap:background} --- Background and Concepts}
presents the theoretical and technical foundations of the work.  It
introduces sequential pattern matching and its applications, explains the
eight \texttt{MATCH\_RECOGNIZE} sub-clauses, and presents the SQL:2016
R010--R030 feature classification used in this thesis.

\textbf{Chapter~\ref{chap:related} --- Related Work}
reviews Oracle, Trino, Snowflake, Apache Flink, and DuckDB.  It compares
their feature scope, execution strategies, deployment models, and practical
restrictions.  It then examines the Python and Pandas gap and positions this
work within the row-pattern literature.

\textbf{Chapter~\ref{chap:methodology} --- System Design and Methodology}
describes the architecture and methodology.  It covers the SQL grammar,
parser, AST, semantic validation, logical plan, and the design that connects
a declarative query to automata-based execution.

\textbf{Chapter~\ref{chap:implementation} --- Implementation and Optimization}
explains NFA and DFA construction, Pandas execution, output construction,
and measure evaluation.  It also presents the performance layer: vectorized
conditions, compiled fast paths, pattern caching, compact output
construction, and memory management.

\textbf{Chapter~\ref{chap:evaluation} --- Experimental Evaluation}
presents the methodology, datasets, and results.  It evaluates correctness,
feature and pattern coverage, usability, pattern complexity, performance,
memory, and scalability.  The chapter compares five pattern families across
six Amazon UK Products sizes and then presents a larger Amazon Reviews
scalability test.  Trino and Oracle provide the cross-system references.

\textbf{Chapter~\ref{chap:conclusions} --- Conclusions and Future Work}
summarizes the contributions and current limitations.  It also identifies
directions for future research and extension.

\section{Publications in the Course of this Thesis}

This thesis produced two scientific outputs in data systems and sequential
pattern recognition: one published paper and one submitted manuscript.

\begin{description}

\item[First Publication.]
The first paper was published in the
\textit{2026 IEEE 23rd Mediterranean Electrotechnical Conference
(MELECON)}.  The paper, entitled ``Row Pattern Recognition for
Pandas: Bringing MATCH\_RECOGNIZE to Python DataFrames''
\cite{lawande2026rpr}, addresses the lack of native row-pattern recognition
in Python data-analysis frameworks.  It presents an SQL-in-Pandas engine for
running \texttt{MATCH\_RECOGNIZE}-style queries directly on DataFrames.  The
paper describes its supported row-pattern features, NFA-based matching method,
and initial tests.  Those tests show that the declarative interface reduces
procedural Pandas code while preserving the expected results in the studied
cases.

\item[Second Publication.]
A second manuscript extending this work has been submitted to
\textit{Scientific Reports} and is currently under review.

\end{description}

% ----------------------------------------------------------------
\section{Chapter Summary}
\label{sec:intro-summary}
% ----------------------------------------------------------------

This chapter introduced row-pattern recognition and the lack of a native,
declarative equivalent in Pandas.  It explained why hand-written multi-row
solutions become verbose, fragile, and difficult to extend.  It then stated
the research objectives, summarized the contributions and evidence, outlined
the thesis, and listed the research outputs.  Chapter~\ref{chap:background}
next introduces the theory and SQL semantics used by the system.
